% Build with LuaLaTeX (or XeLaTeX), twice:   lualatex example.tex
%
% For a question-only copy, change  solutions=show  to  solutions=hide
% below (or load the package with the `nosolutions' option).
\documentclass[11pt]{article}

\usepackage[a4paper, left=14mm, right=24mm, top=16mm, bottom=22mm]{geometry}
\usepackage{fontspec}
\setmainfont{TeX Gyre Pagella}
\usepackage{amsmath, amssymb}
\usepackage{passign}   % after geometry

\passignsetup{
  course        = Linear Algebra,
  course code   = MATH 2210,
  term          = Autumn 2026,
  number        = 3,
  title         = {Eigenvalues and\\ Diagonalisation},
  author        = Ada Lovelace,
  student id    = 20260042,
  instructor    = Prof.\ C.\ Babbage,
  due           = 12 October 2026,
  solutions     = show,
}

\newcommand{\R}{\mathbb{R}}

\begin{document}

\maketitle
\addtitlefield{Collaborators}{None}

\begin{remark*}[Conventions]
Throughout, matrices are real unless stated otherwise, and $I$ denotes
the identity matrix of the appropriate size.  Show all working.
\end{remark*}

\section{Warm-up}

\begin{question}[title=A first eigenvalue computation, marks=10]
Let
\[
  A = \begin{pmatrix} 2 & 1 \\ 1 & 2 \end{pmatrix}.
\]
\begin{parts}
  \part[4] Find the eigenvalues of $A$.
  \part[3] Find a basis of each eigenspace.
  \part[3] Is $A$ diagonalisable?  Justify your answer.
    \begin{subparts}
      \subpart[1] State the criterion you are using.
      \subpart[2] Apply it to $A$.
    \end{subparts}
\end{parts}
\end{question}

\begin{solution}
\begin{parts}
  \part The characteristic polynomial is
  \[
    \det(A-\lambda I) = (2-\lambda)^2 - 1 = (\lambda-1)(\lambda-3),
  \]
  so the eigenvalues are $\lambda_1 = 1$ and $\lambda_2 = 3$.

  \part For $\lambda_1 = 1$ we solve $(A-I)v=0$, which gives $v_1 =
  (1,-1)^{\top}$.  For $\lambda_2 = 3$ we solve $(A-3I)v=0$, which gives
  $v_2 = (1,1)^{\top}$.

  \part Yes.
    \begin{subparts}
      \subpart A matrix is diagonalisable if and only if $\R^n$ has a
      basis of eigenvectors.
      \subpart The vectors $v_1, v_2$ are linearly independent, so they
      form a basis of $\R^2$ and
      \[
        A = PDP^{-1}, \quad
        P = \begin{pmatrix} 1 & 1 \\ -1 & 1 \end{pmatrix}, \quad
        D = \begin{pmatrix} 1 & 0 \\ 0 & 3 \end{pmatrix}.
      \]
    \end{subparts}
\end{parts}
\end{solution}

\section{Proofs}

\begin{question}[5]
Recall the following definition.

\begin{definition}[Eigenvalue]
A scalar $\lambda$ is an \emph{eigenvalue} of $A$ if $Av=\lambda v$ for
some non-zero vector $v$.
\end{definition}

Show that if $\lambda$ is an eigenvalue of an invertible matrix $A$, then
$\lambda \neq 0$ and $\lambda^{-1}$ is an eigenvalue of $A^{-1}$.

\begin{hint*}
Apply $A^{-1}$ to both sides of $Av = \lambda v$.
\end{hint*}
\end{question}

\begin{solution}
Let $Av=\lambda v$ with $v\neq 0$.  If $\lambda=0$ then $Av=0$, which
contradicts the invertibility of $A$.  Hence $\lambda\neq0$, and applying
$A^{-1}$ gives $v=\lambda A^{-1}v$, that is
$A^{-1}v=\lambda^{-1}v$.\marginnote{We need $v\ne0$ here, otherwise the
conclusion is vacuous.}

\begin{remark}
The converse also holds: apply the same argument to $A^{-1}$.
\end{remark}
\end{solution}

\begin{question}[title=Trace and determinant, marks=3+3]
Let $A$ be a $2\times2$ matrix with eigenvalues $\lambda_1,\lambda_2$.
\begin{parts}
  \part Prove that $\operatorname{tr}A=\lambda_1+\lambda_2$.
  \part Prove that $\det A=\lambda_1\lambda_2$.
\end{parts}
\end{question}

\begin{solution}[Via the characteristic polynomial]
Write $\chi_A(\lambda)=\lambda^2-(\operatorname{tr}A)\lambda+\det A$.
Since $\lambda_1,\lambda_2$ are its roots,
$\chi_A(\lambda)=(\lambda-\lambda_1)(\lambda-\lambda_2)$, and comparing
coefficients yields both claims.\sidenote{Valid over $\mathbb{C}$, where the
polynomial always splits.}
\end{solution}

\begin{question*}[label=Bonus, marks=2]
Does every real $2\times2$ matrix have a real eigenvalue?  Give a proof
or a counterexample.
\end{question*}

\begin{solution}
No.  The rotation $R=\begin{pmatrix}0&-1\\1&0\end{pmatrix}$ has
characteristic polynomial $\lambda^2+1$, which has no real roots.
\end{solution}

\end{document}
